\documentclass[11pt,a4paper]{article}

% ─── Packages ────────────────────────────────────────────────────────────────
\usepackage[T1]{fontenc}
\usepackage{lmodern}
\usepackage[margin=2.5cm]{geometry}
\usepackage[pdftex,colorlinks=true,linkcolor=blue!60!black,
            urlcolor=blue!60!black,pdfusetitle]{hyperref}
\usepackage{microtype}
\usepackage{enumitem}
\usepackage{xcolor}
\usepackage{mdframed}
\usepackage{fancyhdr}

\pagestyle{fancy}
\fancyhf{}
\fancyhead[C]{\small\textit{IFT GPU Workstation --- Internal Technical Memo}}
\fancyfoot[C]{\thepage}
\renewcommand{\headrulewidth}{0.4pt}

\setlist[itemize]{leftmargin=*}
\setlist[enumerate]{leftmargin=*}

% ─── Highlight boxes ──────────────────────────────────────────────────────────
\definecolor{boxbg}{RGB}{240,246,255}
\definecolor{boxframe}{RGB}{70,130,180}

\newenvironment{keybox}[1]{%
  \begin{mdframed}[backgroundcolor=boxbg,linecolor=boxframe,linewidth=1pt,
    innertopmargin=6pt,innerbottommargin=6pt,
    innerrightmargin=8pt,innerleftmargin=8pt]
  \textbf{#1.}\quad}{\end{mdframed}}

% ─── Metadata ─────────────────────────────────────────────────────────────────
\title{\Large\bfseries IFT GPU Workstation Memo}
\author{CQMP}
\date{\today}

% ═════════════════════════════════════════════════════════════════════════════
\begin{document}

\maketitle

% ─────────────────────────────────────────────────────────────────────────────
\section{Hardware Configuration}
\label{sec:hardware}
% ─────────────────────────────────────────────────────────────────────────────

The workstations are Dell Pro Max Tower T2 systems installed from the shared
workstation image.  The currently installed reference machine, Matsubara, has
the following configuration:

\begin{itemize}
  \item CPU: Intel Core Ultra 7 265, 20 cores, one socket.
  \item Memory: 30 GiB RAM plus 8 GiB compressed swap.
  \item Local storage: one Kioxia XG10d 1 TB NVMe device.  The standard
        installation uses a 512 MB EFI partition, a 1 GB \texttt{/boot}
        partition, a 50 GB operating-system partition, and the remainder as
        \texttt{/data}.
  \item Graphics: integrated Intel Arrow Lake-S graphics plus an NVIDIA
        RTX 2000 Ada Generation GPU.
  \item Network: Intel I219-LM Ethernet.
\end{itemize}

The workstation inventory and assigned addresses are tracked in
\texttt{machines.md} in the image repository.

% ─────────────────────────────────────────────────────────────────────────────
\section{System Configuration}
\label{sec:system}
% ─────────────────────────────────────────────────────────────────────────────

The machines run CentOS Stream 9 as a bootc image.  In this model, the
operating system is built as a container image and deployed atomically to the
host.  The machine normally boots the current deployment and keeps a rollback
deployment available.  Updates are delivered by publishing a new image and then
running a bootc upgrade on the workstation.

\begin{keybox}{Image repository}
\url{https://github.com/CQMP/workstation-image}
\end{keybox}

\begin{keybox}{Published image}
\texttt{ghcr.io/cqmp/centos9-workstation:latest}
\end{keybox}

The main image definition is the \texttt{Containerfile}.  The unattended
installer is described by \texttt{ks.cfg}; it partitions the local NVMe disk,
pulls the image above, and installs the bootc system.  The image includes GNOME
and GDM, NVIDIA drivers and CUDA tooling, SSH, LDAP/SSSD authentication,
autofs/NFS integration, HTCondor, compilers, MPI stacks, Python scientific
packages, browsers, editors, and common desktop tools.

The image also carries kernel arguments and initramfs configuration needed for
these systems.  In particular, nouveau is blacklisted, NVIDIA DRM mode setting
is enabled, and the initramfs is kept small enough for reliable EFI boot.

% ─────────────────────────────────────────────────────────────────────────────
\section{Package Requests and Image Changes}
\label{sec:changes}
% ─────────────────────────────────────────────────────────────────────────────

Users should not install long-lived system packages manually on individual
workstations.  Manual changes are local to one machine and may disappear after
reinstallation or diverge from the rest of the pool.  Instead, package requests
and system configuration changes should be made through the image repository.

The preferred workflow is:

\begin{enumerate}
  \item Fork \url{https://github.com/CQMP/workstation-image}.
  \item Modify the \texttt{Containerfile} or the relevant file under
        \texttt{etc/}, \texttt{usr/}, or another tracked directory.
  \item Build or inspect the change locally when practical.
  \item Submit a pull request against the main repository.
\end{enumerate}

After a pull request is merged, GitHub Actions rebuilds and publishes the image.
Machines can then receive the update with:

\begin{verbatim}
sudo bootc upgrade
sudo reboot
\end{verbatim}

% ─────────────────────────────────────────────────────────────────────────────
\section{Logging In}
\label{sec:login}
% ─────────────────────────────────────────────────────────────────────────────

Authentication is handled through the FUW/ITP LDAP system via SSSD.  Users log
in with their Unix password.  Being present in LDAP is necessary but not
sufficient: the image also contains an explicit allow list of users who may log
in to these workstations.

The allow list is configured in \texttt{etc/sssd/sssd.conf} as
\texttt{simple\_allow\_users}.  Users who need access should provide their Unix
user name so it can be added to the image.

\begin{keybox}{Current image allow list}
\texttt{egull, rfarid, ajazdzewska, amarie, abalbi}
\end{keybox}

Home directories and shared file systems are resolved through LDAP/autofs.  If
login fails, the first things to check are whether the user exists in the LDAP
Unix domain, whether the user is on the image allow list, and whether the user
is using the Unix password rather than another institutional password.

% ─────────────────────────────────────────────────────────────────────────────
\section{Condor}
\label{sec:condor}
% ─────────────────────────────────────────────────────────────────────────────

Each workstation is configured as an HTCondor execute node.  The pool is managed
by a central manager at \texttt{condor.gull-group.org}.  All ten workstations
register with it using CCB so that the manager can reach them through the campus
firewall.  Each machine runs only the \texttt{MASTER} and \texttt{STARTD}
daemons; the central manager handles scheduling.  Jobs may be submitted from any
workstation in the pool.

\subsection{Checking pool status}

\begin{verbatim}
condor_status                         # one line per slot
condor_status -constraint 'GPUs >= 1' # GPU slots only
\end{verbatim}

The \texttt{Activity} column shows \texttt{Idle} (ready to accept work),
\texttt{Busy} (running a job), \texttt{Suspended} (job paused, local user
active), or \texttt{Vacating} (checkpointing before eviction).

\subsection{Writing a submit file}

A minimal CPU submit file:

\begin{verbatim}
executable = /bin/hostname
log        = job.log
output     = job.out
error      = job.err
queue
\end{verbatim}

For a GPU job, add resource requests:

\begin{verbatim}
executable     = my_cuda_program
log            = job.log
output         = job.out
error          = job.err
request_gpus   = 1
request_cpus   = 4
request_memory = 8 GB
queue
\end{verbatim}

Key directives:

\begin{description}
  \item[\texttt{executable}] Path to the binary or script.  Relative paths are
        resolved from the submit directory.
  \item[\texttt{arguments}] Command-line arguments, space-separated.
  \item[\texttt{log}] HTCondor event log: submit (\texttt{000}), execute
        (\texttt{001}), evict (\texttt{004}), hold (\texttt{009}), and
        normal termination (\texttt{005}) events all appear here.
  \item[\texttt{output}, \texttt{error}] Where the job's stdout and stderr
        are written after the job completes (or are streamed back if
        \texttt{stream\_output = true} is set).
  \item[\texttt{request\_gpus}] Number of GPUs required.  HTCondor sets
        \texttt{CUDA\_VISIBLE\_DEVICES} so the job sees only its assigned device.
  \item[\texttt{request\_memory}] RAM reservation; the job is held if it
        substantially exceeds this.
  \item[\texttt{transfer\_input\_files}] Comma-separated list of files to copy
        to the execute node before the job starts.
  \item[\texttt{transfer\_output\_files}] Files to copy back when the job
        finishes.  If omitted, all new files in the job sandbox are returned.
\end{description}

Files listed relative to the submit file are transferred automatically.  For
large datasets it is more efficient to reference the shared NFS mounts
(\texttt{/dmj}, \texttt{/expo}) directly from the job, or to pre-stage data
on the execute node's local \texttt{/data} scratch disk.

\subsection{Submitting and monitoring jobs}

\begin{verbatim}
condor_submit job.sub           # submit; prints cluster.proc ID
condor_q                        # list your queued and running jobs
condor_q -better-analyze        # explain why a job is not starting
condor_watch                    # live queue view (refreshes every 2 s)
condor_rm  <cluster>            # remove a job or all jobs in a cluster
condor_hold <cluster>           # pause without removing
condor_release <cluster>        # resume a held job
\end{verbatim}

The \texttt{log} file is the primary diagnostic tool: it records every state
transition with a timestamp.  If a job is evicted and restarted the log shows
both the eviction event (\texttt{004}) and the subsequent execute event
(\texttt{001}) on the new machine.

\subsection{Machine state and eviction}
\label{sec:eviction}

The workstations are shared with interactive users.  The execute-node config
defines the following transitions:

\begin{keybox}{Start / suspend / evict policy}
\small\setlength{\parskip}{2pt}
\begin{description}[leftmargin=2.2cm,style=nextline,font=\normalfont\ttfamily]
  \item[START]    KeyboardIdle $>$ 900\,s \& LoadAvg $<$ 0.5 --- job accepted
  \item[SUSPEND]  KeyboardIdle $<$ 120\,s --- SIGSTOP, job frozen in place
  \item[CONTINUE] KeyboardIdle $>$ 600\,s \& LoadAvg $<$ 0.5 --- SIGCONT, resumes
  \item[PREEMPT]  Active $>$60\,s after suspend --- checkpoint initiated, then evict
  \item[KILL]     After 1\,h vacate time --- SIGKILL if still running
\end{description}
\end{keybox}

The two-stage design avoids unnecessary evictions.  When a local user returns,
Condor first \emph{suspends} the job (cheap: the process is frozen in place,
the GPU sits idle, nothing is written to disk).  If the user leaves again
within about ten minutes the job is simply resumed with no overhead.  Only if
the user stays active does Condor escalate to \emph{preemption}: it saves a
checkpoint (described below), evicts the job, and requeues it.  The job will
restart from the checkpoint on any available slot in the pool.

\subsection{CRIU checkpointing and GPU state preservation}
\label{sec:checkpointing}

\emph{Transparent checkpointing} lets a running process be frozen, written to
disk, transferred to a different machine, and resumed --- without any changes to
the application source code.  The mechanism has two layers:

\begin{itemize}
  \item \textbf{CRIU} (Checkpoint/Restore In Userspace) captures the full
        CPU-side process state: memory pages, file descriptors, open sockets,
        threads, and signal masks.
  \item \textbf{\texttt{cuda-checkpoint}} (NVIDIA) quiesces the GPU context,
        saves device registers and VRAM allocations, and restores them after
        CRIU restores the CPU process on the destination node.
\end{itemize}

Both tools are installed by the workstation image.
\texttt{ENABLE\_CHECKPOINTING = True} in the execute-node config allows the
\texttt{condor\_starter} to invoke them.  To opt a job into transparent
checkpointing, add one line to the submit file:

\begin{verbatim}
want_checkpointing = true
\end{verbatim}

When eviction is triggered, HTCondor:
\begin{enumerate}
  \item Calls \texttt{cuda-checkpoint -{}-action checkpoint} to freeze the GPU
        context and flush it to host memory.
  \item Runs CRIU \texttt{dump} to snapshot the process (CPU state + the GPU
        image just flushed) and write it to the job sandbox on local NVMe.
  \item Transfers the checkpoint bundle to the schedd spool directory on
        \texttt{condor.gull-group.org}, so the checkpoint survives if the execute
        node is rebooted or reimaged.
  \item Requeues the job.  When a matching slot becomes available, the bundle is
        transferred there and CRIU \texttt{restore} + \texttt{cuda-checkpoint
        -{}-action restore} restart the job from the exact point it was frozen.
\end{enumerate}

The RTX~2000~Ada has 8~GB of VRAM.  At the NVMe write bandwidth of these
machines ($\sim$3~GB/s) a fully-loaded GPU takes a few seconds to checkpoint to
local disk; the one-hour \texttt{MachineMaxVacateTime} gives ample margin before
the forced kill.

\begin{keybox}{Checkpointing limitations}
Not all CUDA workloads can be checkpointed.  The following are known to prevent
a successful \texttt{cuda-checkpoint}: NCCL collective operations in flight,
\texttt{cudaIpcMemHandle} shared-memory handles, and some third-party GPU
libraries that manage their own CUDA contexts.  If \texttt{cuda-checkpoint}
fails, HTCondor falls back to killing the job and restarting it from scratch on
the next available slot.
\end{keybox}

\subsection{Example jobs}
\label{sec:examples}

Complete submit files and source code are in \texttt{examples/condor/} in the
image repository.

\subsubsection*{Example~1: CPU sanity check}

Verifies the pool is reachable and basic scheduling works.

\begin{verbatim}
# examples/condor/01-cpu-sanity.sub
executable = /bin/hostname
log        = 01-cpu.log
output     = 01-cpu.out
error      = 01-cpu.err
queue
\end{verbatim}

\begin{verbatim}
condor_submit 01-cpu-sanity.sub
condor_q
cat 01-cpu.out      # should print a workstation hostname
\end{verbatim}

\subsubsection*{Example~2: GPU slot assignment}

Requests one GPU and runs \texttt{nvidia-smi} to confirm the device is
accessible.  HTCondor sets \texttt{CUDA\_VISIBLE\_DEVICES} so the job sees
only its assigned GPU.

\begin{verbatim}
# examples/condor/02-gpu-info.sub
executable   = /usr/bin/nvidia-smi
log          = 02-gpu.log
output       = 02-gpu.out
error        = 02-gpu.err
request_gpus = 1
queue
\end{verbatim}

The output file will contain the full \texttt{nvidia-smi} table for the
assigned GPU.

\subsubsection*{Example~3: Long-running CUDA job with checkpointing}

\texttt{vec\_add.cu} allocates two 256~MB float arrays on the device and runs
1200~rounds of SAXPY at half-second intervals, printing a timestamped progress
line each round.  Total wall time is approximately ten minutes --- long enough
to trigger a real checkpoint if a user returns to the machine.

Compile on any workstation:

\begin{verbatim}
make -C examples/condor
\end{verbatim}

or directly:

\begin{verbatim}
nvcc -O2 -arch=sm_89 -o vec_add vec_add.cu
\end{verbatim}

Submit file:

\begin{verbatim}
# examples/condor/03-vec-add.sub
executable         = vec_add
log                = 03-vec-add.log
output             = 03-vec-add.out
error              = 03-vec-add.err
request_gpus       = 1
request_cpus       = 1
request_memory     = 1 GB
want_checkpointing = true
queue
\end{verbatim}

Submit and follow progress:

\begin{verbatim}
condor_submit 03-vec-add.sub
tail -f 03-vec-add.out
\end{verbatim}

To test checkpointing explicitly, hold and immediately release the job while
it is running:

\begin{verbatim}
condor_hold <cluster>
condor_release <cluster>
\end{verbatim}

Alternatively, sit down at the execute node and move the mouse for a minute.
Either way, once the job restarts, the round counter in the output will continue
from the round at which the job was frozen --- confirming that the CRIU
checkpoint was successfully saved, transferred, and restored rather than the job
restarting from round~1.

% ─────────────────────────────────────────────────────────────────────────────
\section{File System}
\label{sec:fs}
% ─────────────────────────────────────────────────────────────────────────────

Authentication and network file-system access are handled through ITP/FUW Unix
credentials.  The image configures SSSD for LDAP and autofs for the shared
mount points \texttt{/dmj}, \texttt{/expo}, and \texttt{/repo}.
These network file systems are appropriate for home directories, shared code,
small inputs, and files that need central management.

Each workstation also has a fast local NVMe scratch disk mounted at
\texttt{/data}.  The standard installation creates per-user directories there
for allowed users.  This disk is not backed up.  It should be used for
simulation scratch data, compilation trees, temporary files, local virtual
environments, and other high-I/O work that does not need backup.  Important
results should be copied back to backed-up network storage.

The local scratch disks are also exported within the workstation subnet and are
available on demand under \texttt{/shared\_data}.  For example,
\texttt{/shared\_data/matsubara} accesses the \texttt{/data} disk on
\texttt{matsubara}, \texttt{/shared\_data/hubbard} accesses the \texttt{/data}
disk on \texttt{hubbard}, and similarly for the other hosts listed in
\texttt{machines.md}.  These mounts are managed by autofs, so a host is mounted
only when its directory is accessed.

The \texttt{/shared\_data/<host>} paths are a convenience for moving or reading
scratch data between workstations; they do not make the data backed up or
centrally managed.  The mounts use soft NFS semantics so that access to an
offline workstation fails rather than indefinitely blocking a user's shell or
job.  Users should therefore treat these paths as best suited for interactive
copying, inspection, and recoverable scratch workflows, not as a replacement for
backed-up storage.

\end{document}
